\begin{tikzpicture}[font=\small, >=Latex]
\tikzstyle{root}=[
    rectangle, rounded corners=2pt, draw=DTBlue, fill=DTLightBlue,
    text width=5.0cm, minimum height=0.9cm, align=center, line width=0.7pt
]
\tikzstyle{cat}=[
    rectangle, rounded corners=2pt, draw=DTGreen, fill=DTLightGreen,
    text width=4.1cm, minimum height=1.65cm, align=left, font=\scriptsize
]

\node[root] (root) at (0,0) {Digital Twin applications in eHealth};

\node[cat] (patient) at (-5.2,-2.0) {\textbf{Patient-specific twins}\\Data: EHR, wearables, labs\\Services: monitoring, risk, personalization};
\node[cat] (organ) at (0,-2.0) {\textbf{Organ and body-system twins}\\Data: imaging, signals, physiology\\Services: simulation, intervention planning};
\node[cat] (pathway) at (5.2,-2.0) {\textbf{Disease and care-pathway twins}\\Data: events, treatments, outcomes\\Services: trajectory analysis, care coordination};
\node[cat] (device) at (-5.2,-4.35) {\textbf{Medical-device twins}\\Data: device telemetry, usage, alarms\\Services: maintenance, safety, optimization};
\node[cat] (hospital) at (0,-4.35) {\textbf{Hospital and workflow twins}\\Data: resources, queues, operations\\Services: scheduling, capacity, process simulation};
\node[cat] (population) at (5.2,-4.35) {\textbf{Public-health and population twins}\\Data: cohorts, environment, mobility\\Services: prevention, policy, emergency response};

\foreach \n in {patient,organ,pathway,device,hospital,population}{
    \draw[->, thick, DTBlue] (root) -- (\n);
}
\end{tikzpicture}
